\documentclass[aps,prd,reprint,preprintnumbers,showpacs,floatfix,nofootinbib,superscript address,longbibliography]{revtex4-1}
\pdfoutput=1

\usepackage[utf8x]{inputenc}
\usepackage{scalerel,stackengine}
\usepackage{titlesec}
\usepackage{parskip}
\usepackage{stix}
\usepackage{amssymb}
\usepackage{amsmath}
\usepackage{amsfonts}
\usepackage{hhline}
\usepackage{mathtools}
\usepackage[dvipsnames]{xcolor}

\usepackage{pifont}
\newcommand{\cmark}{\hspace{5pt}\color{black}\ding{51}\hspace{5pt}}%
\newcommand{\xmark}{\hspace{5pt}\color{black}\ding{55}\hspace{5pt}}%
\newcommand{\UnitaryTensor}{\color{SeaGreen}$\mathbf{2}$\color{black}}%
\newcommand{\UnitaryVector}{\color{SeaGreen}$\mathbf{3}$\color{black}}%
\newcommand{\GhostVector}{\color{red}$\mathbf{3_{\mathrm{\mathbf{g}}}}$\color{black}}%
\newcommand{\UnitaryStrongVector}{\color{red}$\mathbf{3}_{\mathrm{\mathbf{s}}}$\color{black}}%
\newcommand{\GhostStrongVector}{\color{red}$\mathbf{3}_{\mathrm{\mathbf{g,s}}}$\color{black}}%
\newcommand{\ArrowForwards}[2]{%
	\parbox{#1}{\tikz{\draw[very thick, ->](0,0) to [out=-90,in=270,looseness=#2] (#1,0);}}
}
\newcommand{\ArrowBackwards}[2]{%
	\parbox{#1}{\tikz{\draw[very thick, <-](0,0) to [out=-90,in=270,looseness=#2] (#1,0);}}
}
\usepackage{xspace}
\usepackage{multirow,tabularx}
\usepackage{siunitx}
\usepackage{multirow}
\usepackage{graphicx}
\usepackage{xstring}
\usepackage{etoolbox}
\usepackage{notoccite}
\usepackage{natbib}
\usepackage{mathrsfs}
\usepackage{lineno}
\usepackage{tensor}
\usepackage{accents}
\usepackage{tikz}

\allowdisplaybreaks

\parskip 1mm
\parindent 2mm

\input{manuscriptmacros.tex} %some personal macros

\usepackage{hyperref}
\hypersetup{%
     colorlinks = true,%
     linkcolor = blue,%
     citecolor = blue,%
     filecolor = blue,%
     urlcolor = blue% 
     }%
\usepackage[capitalize]{cleveref} %always load this last in preamble

\begin{document}

\title{Einstein-Proca theory from the Einstein-Cartan formulation}

\author{Will Barker}
\email{wb263@cam.ac.uk}
\affiliation{Astrophysics Group, Cavendish Laboratory, JJ Thomson Avenue, Cambridge CB3 0HE, UK}
\affiliation{Kavli Institute for Cosmology, Madingley Road, Cambridge CB3 0HA, UK}
\author{Sebastian Zell}
\email{sebastian.zell@uclouvain.be}
\affiliation{Centre for Cosmology, Particle Physics and Phenomenology -- CP3,
	Universit\'e catholique de Louvain, B-1348 Louvain-la-Neuve, Belgium}

%\date{}

\begin{abstract}
	We construct a theory of gravity in which a propagating massive vector field arises from a quadratic curvature invariant. 
	The Einstein-Cartan formulation and a partial suppression of torsion ensure the absence of ghost and strong-coupling problems, as we prove with nonlinear Lagrangian and Hamiltonian analysis. Augmenting General Relativity with a propagating torsion vector, our theory provides a purely gravitational origin of Einstein-Proca models and constrains their parameter space. As an outlook to phenomenology, we discuss the gravitational production of fermionic dark matter.
\end{abstract}



\maketitle

\section{Open questions in General Relativity}
The theory of General Relativity (GR) provides an outstandingly successful description of gravity and has been confirmed by countless experiments, such as the breakthrough discovery of gravitational waves~\cite{LIGOScientific:2016aoc}. Nevertheless, many far-reaching questions have remained unanswered. Gravitational interactions have long revealed the existence of a non-baryonic form of matter~\cite{Zwicky:1933gu}. It is unknown, however, what this dark matter is composed of and how it was produced in the cosmological history~\cite{Feng:2010gw,Peebles:2013hla,Bertone:2016nfn,Arbey:2021gdg}. Moreover, the very early moments of our Universe -- due to a suspected phase of inflation~\cite{Starobinsky:1980te, Guth:1980zm, Linde:1981mu, Mukhanov:1981xt} -- and its far future -- because of the observed dark energy~\cite{SupernovaSearchTeam:1998fmf, SupernovaCosmologyProject:1998vns} -- appear to have in common an accelerated expansion. Agreement with recent	measurements, in particular of the cosmic microwave background (CMB)~\cite{Planck:2018jri, BICEP:2021xfz}, is excellent but a microscopic origin has still not been determined both for dark energy and for inflation.

It is exciting to explore if GR also contains an answer to these problems.  As we shall show, this question is closely connected to the geometry of spacetime, which a priori is characterized by three independent properties curvature $\rcR{}$, torsion $T^\alpha_{~~\beta \gamma}$ and nonmetricity $Q_{\alpha\mu\nu}$ (see~\cref{fig:schema}). The presence or absence of $\rcR{}$, $Q_{\alpha\mu\nu}$ and $T^\alpha_{~~\beta \gamma}$ defines seven formulations of GR~\cite{Einstein:1915,Weyl:1918, Palatini:1919,Weyl:1922,Cartan:1922, Eddington:1923,  Cartan:1923,Eddington:1923, Cartan:1924, Cartan:1925,Einstein:1925, Einstein:1928, Einstein:19282,Schroedinger:1950,Moller:1961, Pellegrini:1963, Hayashi:1967se,Cho:1975dh,Hehl:1976kt, Hehl:1976kv, Hehl:1976my, Hehl:1977fj,Kijowski:1978,Hayashi:1979qx, Nester:1998mp, BeltranJimenez:2019odq} (see~\cite{Heisenberg:2018vsk, BeltranJimenez:2019esp, Rigouzzo:2022yan,Heisenberg:2023lru} for an overview).  
For example, assuming that both $Q_{\alpha\mu\nu}$ and $T^\alpha_{~~\beta \gamma}$ vanish leads to the commonly-used metric variant of GR~\cite{Einstein:1915} corresponding to Riemannian geometry, whereas Einstein-Cartan (EC) gravity is defined by including both $\rcR{}$ and $T^\alpha_{~~\beta \gamma}$ while excluding $Q_{\alpha\mu\nu}$~\cite{Cartan:1922,Cartan:1923,Cartan:1924,Cartan:1925,Einstein:1925, Einstein:1928,Einstein:19282}. The EC version stands out since it can be derived by gauging the Poincar\'e group~\cite{Utiyama:1956sy, Kibble:1961ba, Sciama:1962}, which puts gravity on the same footing as the other fundamental forces.
This fact -- together with its phenomenological virtues -- motivates us to focus on EC gravity in the following.
 
	\begin{figure}
		\vspace{-0.5\baselineskip}
	\includegraphics[width=1\linewidth]{NonRiemannianSchematic.pdf}
		\vspace{-2.5\baselineskip}
		\caption{Schematic representation of parallel transport in the presence of curvature $\tensor{R}{^\rho_{\sigma\mu\nu}}$ -- rotation around closed curves -- torsion $\tensor{T}{^\sigma_{\mu\nu}}$ -- non-closure of (infinitesimal) parallelograms -- and nonmetricity $\tensor{Q}{_{\lambda\mu\nu}}$ -- non-conservation of vector norms.}
	\label{fig:schema}
\end{figure}

 At first sight, the various versions of GR appear to be very different. In the classical, matter-free limit, however, they are all fully equivalent.
There are two main ways to break this equivalence. First, one can couple matter (non-minimally) to GR. In this case, the spectrum of gravity remains unchanged, \ie no additional propagating \dof emerge apart from the massless spin-2 graviton, but new effects arise at high energies~\cite{Freidel:2005sn,Bauer:2008zj,Poplawski:2011xf,Diakonov:2011fs,Khriplovich:2012xg,Magueijo:2012ug,Khriplovich:2013tqa,Markkanen:2017tun,Carrilho:2018ffi,Enckell:2018hmo,Rasanen:2018fom,Rubio:2019ypq,Shaposhnikov:2020geh,Karananas:2020qkp,Langvik:2020nrs,Shaposhnikov:2020gts,Mikura:2020qhc,Shaposhnikov:2020aen,Kubota:2020ehu,Enckell:2020lvn,Iosifidis:2021iuw,Racioppi:2021ynx,Cheong:2021kyc,Dioguardi:2021fmr,Piani:2022gon,Dux:2022kuk,Rasanen:2022ijc,Gialamas:2023emn,Gialamas:2023flv,Piani:2023aof,Poisson:2023tja,Rigouzzo:2023sbb}. An example, which we shall discuss later, is as a novel mechanism for producing dark matter in EC gravity~\cite{Shaposhnikov:2020aen}.  Second, one can add terms nonlinear in curvature. Almost inevitably, this leads to the existence of new dynamical particles~\cite{Stelle:1977ry,Neville:1978bk,Neville:1979rb,Sezgin:1979zf,Hayashi:1979wj, Hayashi:1980qp}. 
These additional propagating \dof are generically plagued by various types of inconsistencies~\cite{Stelle:1977ry,Neville:1978bk,Neville:1979rb,Sezgin:1979zf,Hayashi:1979wj, Hayashi:1980qp,Hecht:1990wn,Hecht:1991jh,Chen:1998ad}. 
Finding a model which is consistent and weakly coupled has turned out to be challenging and, for torsion, so far was only achieved for a propagating scalar mode~\cite{Hecht:1996np,Chen:1998ad,Yo:1999ex,Yo:2001sy}. 

In this paper, we will use the concept of partial Lagrange multipliers~\cite{Barker:2022jsh} to obtain for the first time a propagating torsion vector from a curvature-squared term in a model for which consistency can be proven. In addition to a proof of concept for a new class of models, our findings can have implications for phenomenology since dynamical (Abelian) massive vector fields have been employed for inflation~\cite{Ford:1989me,Golovnev:2008cf,Chiba:2008eh,Koivisto:2008xf,Kanno:2008gn,Emami:2016ldl,Rodriguez:2017wkg}, dark energy~\cite{Armendariz-Picon:2004say,Koivisto:2007bp,Koivisto:2008xf,Tasinato:2014eka,DeFelice:2016yws,DeFelice:2016uil,BeltranJimenez:2016afo,deFelice:2017paw,Rodriguez:2017wkg,Nakamura:2018oyy,Heisenberg:2020xak,Benisty:2021sul,deRham:2021efp} and dark matter~\cite{Hambye:2008bq,Hambye:2009fg,Arina:2009uq,Hisano:2010yh,Diaz-Cruz:2010czr,Lebedev:2011iq,Farzan:2012hh,Baek:2012se,Belyaev:2016icc,Arcadi:2020jqf,Barman:2021qds}, among others (see also~\cite{Bekenstein:1971hc,Heisenberg:2017hwb,Garcia-Saenz:2021uyv} for implications for black holes). In these works, however, an a priori independent vector field was considered as an ad hoc addition to GR. Our theory endows these Einstein-Proca models with a gravitational origin. This is not only conceptually appealing but also leads to constraints on parameters, which would otherwise seem arbitrary from an effective field theory perspective. Moreover, we show that the massive vector mode can be integrated out in relevant parts of parameter space, thereby establishing a connection to works with non-propagating torsion using the example of~\cite{Shaposhnikov:2020aen}.

\section{Obstructions to dynamical torsion}
The simplest choice of action for EC gravity is (see~\cite{Hehl:1976kj, Shapiro:2001rz, Blagojevic:2002, Blagojevic:2012bc} for reviews, and \cref{Conventions} for details of our conventions)
\begin{equation} \label{EC}
	L_{\text{EC}}  \equiv -\frac{1}{2} \Planck^2 \rcR{} 	+ \TLambda[Q]{_{\mu\nu\sigma}}Q{^{\mu\nu\sigma}}+L_{\text{M}}\left(\Gamma\right) \;,
\end{equation}
where the multiplier field $\TLambda[Q]{_{\mu\nu\sigma}}$ enforces the vanishing of nonmetricity and we included matter $L_{\text{M}}\left(\Gamma\right)$. We can decompose curvature in its Riemannian part $\rR{}$ and contributions from irreducible parts of torsion:
	\begin{align}
		\rcR{}\left(Q_{\alpha\mu\nu}\mapsto 0\right) & = \rR{}
		+
		\frac{8}{9}
		\T[1]{_{\mu\nu\sigma}}
		\T[1]{^{\mu[\nu\sigma]}}
		-
		\frac{2}{3}
		\T[2]{_{\mu}}
		\T[2]{^{\mu}}	
		\nonumber\\
		&\ \ \ 
		+
		\frac{3}{2}
		\T[3]{_{\mu}}
		\T[3]{^{\mu}}	
		-
		2\rD{_\mu}\T[2]{^\mu} \;.
\label{decomposition}
	\end{align}
Varying~\cref{EC} \wrt torsion and without matter, it follows that $\T[1]{^{\mu\nu\sigma}}\approx\T[2]{^{\mu}}\approx\T[3]{^{\mu}}\approx 0$. Plugging this result back in the action shows that the theory \eqref{EC} is equivalent to the metric formulation of GR, $L_{\text{GR}} \equiv -\frac{1}{2} \Planck^2 \rR{}$. 

Adding terms that are quadratic in curvature~\cite{Utiyama:1956sy,Sciama:1964wt,Hayashi:1980ir,Hayashi:1980qp} to~\eqref{EC} generically leads to two types of problems caused by additional propagating d.o.f. The first one consists of ghost and tachyonic instabilities, \ie kinetic or mass terms with the wrong sign~\cite{Stelle:1977ry,Neville:1978bk,Neville:1979rb,Sezgin:1979zf,Hayashi:1979wj, Hayashi:1980qp}. Superficially, one can eliminate these inconsistencies by tuning the kernel of the linearised wave operator, and indeed a `zoo' of such unitary cases is available~\cite{Sezgin:1981xs,Blagojevic:1983zz,Blagojevic:1986dm,Kuhfuss:1986rb,Yo:1999ex,Yo:2001sy,Blagojevic:2002,Puetzfeld:2004yg,Yo:2006qs,Shie:2008ms,Nair:2008yh,Nikiforova:2009qr,Chen:2009at,Ni:2009fg,Baekler:2010fr,Ho:2011qn,Ho:2011xf,Ong:2013qja,Puetzfeld:2014sja,Karananas:2014pxa,Ni:2015poa,Ho:2015ulu,Karananas:2016ltn,Obukhov:2017pxa,Blagojevic:2017ssv,Blagojevic:2018dpz,Tseng:2018feo,Lin:2018awc,BeltranJimenez:2019acz,Zhang:2019mhd,Aoki:2019rvi,Zhang:2019xek,Jimenez:2019qjc,Lin:2019ugq,Percacci:2019hxn,Barker:2020gcp,BeltranJimenez:2020sqf,Barker:2021oez,MaldonadoTorralba:2020mbh,Marzo:2021esg,Marzo:2021iok,delaCruzDombriz:2021nrg,Baldazzi:2021kaf,Annala:2022gtl}. However, a second class of issues, first discussed in the metric-compatible teleparallel formulation of GR~\cite{Einstein:1928, Einstein:19282, Moller:1961, Pellegrini:1963, Hayashi:1967se,Cho:1975dh, Hayashi:1979qx} and its generalizations \cite{Hayashi:1979qx,Dimakis:1989az,Dimakis:1989ba,Lemke:1990su,Hecht:1990wn,Hecht:1991jh,Afshordi:2006ad,Magueijo:2008sx}, is only visible in a nonlinear analysis~\cite{Velo:1969txo,Aragone:1971kh,Cheng:1988zg,Chen:1998ad,Chen:1998ad,Blixt:2018znp,Blixt:2019ene,Blixt:2020ekl,Krasnov:2021zen,Bahamonde:2021gfp,Delhom:2022vae}. Applied to EC gravity, an especially restrictive report~\cite{Yo:2001sy} found that for propagating tensor/vector torsion, weak-field linearisation causes over-counting of the constraints~\cite{Blagojevic:1983zz,Blagojevic:1986dm,Blagojevic:2002,Ong:2013qja,Blagojevic:2018dpz,Barker:2021oez} -- peculiar to torsion theories, the activated fields are ghosts whenever the linear spectrum is unitary~\cite{Hayashi:1980qp,Blagojevic:1983zz,Yo:2001sy,Blagojevic:2013dea,Blagojevic:2013taa,BeltranJimenez:2019hrm,Aoki:2020rae,Barker:2021oez}. In the modern parlance (see~\cref{SCIntro},~\cite{Yo:2001sy,Charmousis:2008ce,Charmousis:2009tc,Papazoglou:2009fj,Baumann:2011dt,Baumann:2011dt,Wang:2017brl,JimenezCano:2021rlu,Barker:2022kdk,Delhom:2022vae,Annala:2022gtl}), this means that these new particles become \emph{strongly coupled} on Minkowski spacetime, which is thereby rendered dynamically unreachable~\cite{DAmico:2011eto,Gumrukcuoglu:2012aa,Mazuet:2017rgq,BeltranJimenez:2020lee,Barker:2022kdk}.\footnote{Such issues of strong coupling have been discovered in a variety of theories~\cite{Vainshtein:1972sx,Deffayet:2001uk,Deffayet:2005ys,Charmousis:2008ce,Charmousis:2009tc,Papazoglou:2009fj,deRham:2014zqa,Deser:2014hga,Wang:2017brl}.} As stated above, models that avoid this problem and remain weakly coupled have so far only been constructed for a propagating \emph{scalar} mode of torsion~\cite{Hecht:1996np,Yo:1999ex,Yo:2001sy} (see e.g.~\cite{Yo:2006qs,Shie:2008ms,Chen:2009at,Baekler:2010fr,Ho:2011qn,Ho:2011xf,Ho:2015ulu,Tseng:2018feo,Zhang:2019mhd,Zhang:2019xek,MaldonadoTorralba:2020mbh,delaCruzDombriz:2021nrg} and~\cite{Puetzfeld:2004yg,Ni:2009fg,Puetzfeld:2014sja,Ni:2015poa,Barker:2020gcp} for applications and reviews, respectively, and~\cite{BeltranJimenez:2019acz,BeltranJimenez:2019esp,Percacci:2020ddy,BeltranJimenez:2020sqf,Marzo:2021esg,Piva:2021nyj,Marzo:2021iok,Iosifidis:2021xdx,Jimenez-Cano:2022sds,Iosifidis:2023pvz} for similar prospects in the nonmetric sector).

In this paper we construct a healthy theory with \emph{propagating vector} (PV) torsion. Extending~\cref{EC}, it is defined
\begin{align}
\LPV{}&\equiv-\frac{1}{2}\Planck^2\rcR{}
+2\alp{5}\rcR{_{[\mu\nu]}}
\rcR{^{[\mu\nu]}}
+\bet{2}\Planck^2\T[2]{_{\mu}}\T[2]{^{\mu}}
\nonumber\\  
	&\hspace{-30pt} 
+\TLambda[1]{_{\mu}^{\nu\sigma}}\T[1]{^{\mu}_{\nu\sigma}}
+\TLambda[3]{^\mu}\T[3]{_\mu}
+\TLambda[Q]{_{\mu\nu\sigma}}Q{^{\mu\nu\sigma}}
+L_{\text{M}}\left(\Gamma\right) \;,\label{DC}
\end{align}
where $\TLambda[1]{_{\mu}^{\nu\sigma}}$ and $\TLambda[3]{^\mu}$ are two Lagrange multipliers and $\alp{5}<0$ throughout.\footnote{One can construct an analogous theory $\LPV{(2)\rightleftarrows(3)}$ with $\alp{5}>0$.} Unlike in recent attempts~\cite{BeltranJimenez:2019acz,BeltranJimenez:2019esp,BeltranJimenez:2020sqf,Jimenez-Cano:2022sds,Iosifidis:2023pvz}, we will demonstrate how consistent dynamics and unitarity are simultaneously upheld by \emph{complete} nonlinear analysis. Our results represent a twofold progress. First, we do not postulate the existence of a kinetic term for $\T[2]{_\mu}$, as is commonly done in EC gravity (see \eg~\cite{Carroll:1994dq,Belyaev:1997zv, Shapiro:2001rz,Vollick:2006uq,Belyaev:2016icc,Barman:2019mlj,Katanaev:2020xdv}). Instead, it arises from the curvature-squared term $\rcR{_{[\mu\nu]}} \rcR{^{[\mu\nu]}}$. Second, it is known that other formulations of GR allow for consistent propagating vectors contained in nonmetricity \cite{BeltranJimenez:2014iie,Iosifidis:2018diy,Iosifidis:2018zwo,Helpin:2019vrv,OrejuelaGarcia:2020viw,Ghilencea:2020piz,BeltranJimenez:2020sih,Xu:2020yeg,Yang:2021fjy,Quiros:2021eju,Quiros:2022uns,Yang:2022icz,Burikham:2023bil,Haghani:2023nrm} (see also~\cite{Aringazin:1991,Vitagliano:2010pq,Iosifidis:2021xdx}) -- indeed our $\rcR{_{[\mu\nu]}}$ in~\cref{DC} is analogous to Weyl's homothetic curvature~\cite{BeltranJimenez:2016wxw,Iosifidis:2018jwu}. Therefore, our theory establishes a direct analogy of \emph{vector torsion} with Weyl's vector nonmetricity.

\section{Nontrivial effect of multipliers} 
Our key innovation in our~\cref{DC} is the \emph{partial} multiplier fields~\cite{Barker:2022jsh}. Without them, the pure $\alp{5}\rcR{_{[\mu\nu]}}\rcR{^{[\mu\nu]}}$ operator propagates $\T[3]{_\mu}$ and $\T[2]{_\mu}$ vectors, both of which are strongly-coupled~\cite{Yo:2001sy,Barker:2022jsh} (absent from the linear spectrum), whilst $\T[3]{_\mu}$ is additionally a ghost (for ${\alpha<0}$). Overlooking strong-coupling, we first naively try to `suppress' the $\T[3]{_\mu}$ ghost solely using the multiplier $\TLambda[3]{^\mu}$ in the theory $L_{\text{PV}}\left(\TLambda[1]{_\mu^{\nu\sigma}},\bet{2}\mapsto 0\right)$. We package six \dof into the 2-form field 
\begin{equation}\label{KalbRamond}
	\begin{aligned}
		\B{_{\mu\nu}}&\equiv 
		\Planck^{-1}\Big(2\rD{_{\sigma}}\T[1]{^{\sigma}_{[\mu\nu]}}
	-\F[2]{_{\mu\nu}}
	+\frac{3}{4}\tensor{\epsilon}{^{\sigma\lambda}_{\mu\nu}}\F[3]{_{\sigma\lambda}}\\
		&\hspace{20pt}
		+2\T[2]{_{\sigma}}\T[1]{^{\sigma}_{[\mu\nu]}}
		+3\tensor{\epsilon}{^{\sigma\lambda\rho}_{[\mu}}\T[3]{_{\sigma}}\T[1]{_{\nu]\lambda\rho}}\Big) \;,
	\end{aligned}
\end{equation}
where the Maxwell field strengths are $\F[2]{_{\mu\nu}}\equiv 2\rD{_{[\mu}}\T[2]{_{\nu]}}$ etc. We introduce~\cref{KalbRamond} just to simplify the $\tensor{\omega}{^{ij}_\mu}$-equations $\delta/\delta \tensor{\omega}{^{ij}_\mu}\int\mathrm{d}^4xeL_{\text{naive}}\approx 0$, which decompose to (see Supplemental Material \cite{FieldEquations} for details about the subsequent computations)
\begin{subequations}
	\begin{align}
		\T[1]{^\sigma_{\mu\nu}}
		&\approx
		\frac{\alp{5}}{\Planck}\left[
		\OmitIndices{\rD{}\B{}}
		+\frac{\OmitIndices{\B{}\rD{}\B{}}}{\Planck}
		+\frac{\OmitIndices{\B{}^2\T[1]{}}}{\Planck}
		\right] \;,\label{TensorEquation}\\
		\T[2]{_\mu}
		&\approx
		\frac{4\alp{5}}{3\Planck}\Big[\rD{_\nu}\B{_\mu^\nu}
		-\B{_{\sigma\lambda}}\T[1]{_{\mu}^{\sigma\lambda}}\Big]\;,\label{VectorEquation}
	\end{align}
\end{subequations}
and the algebraic relation $\TLambda[3]{^\mu}\approx\dots$, with $\T[3]{_\mu}\approx 0$,
where $\OmitIndices{...}$ suppresses contractions, but all parts are simplified by~\eqref{KalbRamond}. With $\alp{5}\mapsto 0$ we recover entirely vanishing vacuum torsion as expected in EC theory, otherwise~\crefrange{TensorEquation}{VectorEquation} should be wavelike for dynamical torsion (if any). It is simplest to notice how all torsion dynamics can be confined to $\B{_{\mu\nu}}$, though that variable eliminates a single derivative in~\cref{KalbRamond}. To extract the propagating (second-derivative) equation in $\B{_{\mu\nu}}$, we take the antisymmetrised divergence of~\eqref{TensorEquation}, next eliminating $\rD{_{\sigma}}\T[1]{^{\sigma}_{[\mu\nu]}}$ for $\B{_{\mu\nu}}$, $\T[1]{^\mu_{\nu\sigma}}$ and $\T[2]{_\mu}$ using~\cref{KalbRamond}, then using~\cref{VectorEquation} to eliminate $\T[2]{_\mu}$ for $\B{_{\mu\nu}}$ and $\T[1]{^\mu_{\nu\sigma}}$, before finally recycling~\cref{TensorEquation} to eliminate all remaining $\T[1]{^\mu_{\nu\sigma}}$ \emph{perturbatively} in terms of $\B{_{\mu\nu}}$. Upon integrating, we find (at least on flat space without matter) that the resulting equation descends from the effective theory 
\begin{align}\label{KalbRamondEffectiveTheory}
	&L_{\text{PV}}\left(\TLambda[1]{_\mu^{\nu\sigma}},\bet{2}\mapsto 0\right)
	\cong 	
	-\frac{\Planck^2}{2}\B{_{\mu\nu}}\B{^{\mu\nu}}
	+2\alp{5}\rD{_{[\mu}}\B{_{\nu\sigma]}}\rD{^{[\mu}}\B{^{\nu\sigma]}}
\nonumber\\
&\hspace{40pt} 
	+\frac{\alp{5}^2}{\Planck}\OmitIndices{\B{}^2\rD{}^2\B{}}
	+\frac{\alp{5}^3}{\Planck^3}\OmitIndices{\B{}^2\rD{}^4\B{}}+\dots
\end{align} 
At~$\mathcal{O}\left(B^2\right)$, \cref{KalbRamondEffectiveTheory} reduces to massive $p$-form electrodynamics~\cite{Henneaux:1986ht} and hence propagates \emph{three} \dof from $\B{_{\mu\nu}}$, while $\rD{^\mu}\big(\delta L_{\text{naive}}/\delta\B{^{\mu\nu}}\big)\propto \rD{_\mu}\B{^\mu_{\nu}}\approx 0$ will constrain the remaining three --- formerly the $\T[3]{_\mu}$ ghost.
Already the fact that the flat-space linear spectrum contains \emph{anything} is remarkable: we accidentally fixed strong coupling of the healthy $\T[2]{_\mu}$ mode by trying to kill the ghost! More strangely still, the ghost is not even dead: the constraint does not survive at~$\mathcal{O}\left(B^3\right)$, so the ghost remains \emph{nonlinearly} active (see~\cref{SCIntro} for a toy model illustrating strong coupling). This conclusion is even true if we ignore the four-derivative term in \cref{KalbRamondEffectiveTheory}, which may lead to additional problems of its own.

\section{Healthy spectrum with multipliers} 
\begin{table}
\begin{tabularx}{\linewidth}{X|c|c|c|c|c|c|c|c}
\hline\hline
	Linear d.o.f & \UnitaryTensor{} & \multicolumn{3}{c|}{\UnitaryTensor{}+\UnitaryVector{}+\GhostVector{}} & \multicolumn{4}{c}{\UnitaryTensor{}+\UnitaryVector{}}\\
\hline
	Nonlinear d.o.f & \UnitaryTensor{}+\UnitaryStrongVector{}+\GhostStrongVector{} & \multicolumn{3}{c|}{\UnitaryTensor{}+\UnitaryVector{}+\GhostVector{}} & \multicolumn{2}{c|}{\UnitaryTensor{}+\UnitaryVector{}+\GhostStrongVector{}} & \multicolumn{2}{c}{\UnitaryTensor{}+\UnitaryVector{}}\\
\hline
	$+\TLambda[3]{^\mu}\T[3]{_\mu}$ &\xmark &\xmark &\xmark &\cmark &\cmark &\xmark &\cmark &\cmark \\
	$+\TLambda[1]{_{\mu}^{\nu\sigma}}\T[1]{^{\mu}_{\nu\sigma}}$ &\xmark &\cmark &\cmark &\xmark &\xmark &\xmark &\cmark &\cmark \\
	$+\bet{2}\Planck^2\T[2]{_{\mu}}\T[2]{^{\mu}}$ &\xmark &\xmark &\cmark &\cmark &\xmark &\cmark &\xmark &\cmark \\
\hline
	Remarks & {\color{red}Failure} & \multicolumn{5}{c|}{{\color{red}Counterintuitive failure}} & \multicolumn{2}{c}{{\color{SeaGreen}Success}}\\
\hline\hline
\end{tabularx}
	\caption{\label{WhacAMole} Additions to $L=-\frac{1}{2}\Planck^2\rcR{}+2\alp{5}\rcR{_{[\mu\nu]}}\rcR{^{[\mu\nu]}}$ needed for a \emph{single} (i.e. non-ghost) $\T[2]{_{\mu}}$ vector in~\cref{DC} with $\alp{5}<0$. Counterintuitively, the pathological $\T[3]{_{\mu}}$ mode \emph{cannot} just trivially be removed via a Lagrange multiplier without activating other ghosts (subscript `g') or strongly-coupled vectors (subscript `s'); see Supplemental Material \cite{NonTrivialMultipliers}.
		All these problems are resolved by our $\TLambda[1]{_{\mu}^{\nu\sigma}}$, which also yields a `conventional' massless limit for $\T[2]{_{\mu}}$ in~\cref{Mass2}.}
\end{table}

Our remarkable experience in~\cref{KalbRamondEffectiveTheory} leads to a highly non-trivial game of `Whac-a-Mole' against strongly-coupled ghosts, whose final score is shown in~\cref{WhacAMole}. 
In another unexpected twist, it turns out that the introduction of the explicit mass parameter $\bet{2}$ has a multiplier-like effect, despite \emph{implicit} $\T[2]{_{\mu}}$ masses being already present within the basic $-\frac{1}{2}\Planck^2\rcR{}$ operator in~\cref{decomposition}. Ultimately, our $\TLambda[1]{_{\mu\nu\sigma}}\T[1]{^{\mu\nu\sigma}}$ term is the cure, even though $\T[1]{^{\mu\nu\sigma}}$ was never implicated in the original pathology~\cite{Yo:2001sy}. For general $\alp{5}$, we restore the whole axial vector sector and, using the spin tensor of matter $e\tensor{S}{^\mu_{ij}}\equiv -2\delta/\delta \tensor{\omega}{^{ij}_\mu}\int\mathrm{d}^4xeL_{\text{M}}$, we obtain the following effective torsion-free theory, to be compared with~\cref{KalbRamondEffectiveTheory}:
\begin{align}\label{MasterLagrangian}
	&\hspace{-15pt}\LPV{\TLambda[3]{^\mu}\mapsto 0}+\bet{3}\Planck^2\T[3]{^\mu}\T[3]{_\mu}
	\nonumber\\
	&
	\cong 
	-\frac{\Planck^2\rR{}}{2}
	+
	\frac{2\alp{5}}{9}
	\F[2]{_{\mu\nu}}
	\F[2]{^{\mu\nu}}
	-
	\frac{\alp{5}}{2}
	\F[3]{_{\mu\nu}}
	\F[3]{^{\mu\nu}}
	\nonumber\\
	&
	+\Planck^2
	\left(
	\tfrac{\left(1+3\bet{2}\right)}{3}
	\T[2]{_{\mu}}
	\T[2]{^{\mu}}	
	-
	\tfrac{\left(3-4\bet{3}\right)}{4}
	\T[3]{_{\mu}}
	\T[3]{^{\mu}}	
	\right)
	\nonumber\\
	&
	-
	\frac{1}{3}
	\T[2]{_{\mu}}
	\Si[2]{^{\mu}}	
	-
	\frac{3}{2}
	\T[3]{_{\mu}}
	\Si[3]{^{\mu}}
	+L_{\text{M}}\left(\rCon{}\right) \;.
\end{align}
In~\cref{MasterLagrangian}, we notice that the residual torsion reduces to the Proca pair, one of which is a ghost, and the full model $\LPV{}$, or $\LPV{(2)\rightleftarrows(3)}$, kills off the ghost in either case. 
In~\cref{DC}, valid for $\alp{5}<0$, the mass of $\T[2]{_\mu}$ is 
\begin{equation}
	\mass[2]{}^2\equiv-3\Planck^2(1+3\bet{2})/4\alp{5} \;.\label{Mass2}
\end{equation}

We shall briefly comment on how our model \eqref{DC} relates to the known formulations of GR. Since some of the irreducible representations of torsion vanish because of partial multipliers~\cite{Barker:2022jsh}, one could classify it as interpolating between the metric and EC formulations, as illustrated in~\cref{Propaganda}. However, our theory \eqref{DC} also shares properties with teleparallel versions of GR, in which curvature is excluded~\cite{Einstein:1928, Einstein:19282, Moller:1961, Pellegrini:1963, Hayashi:1967se,Cho:1975dh, Hayashi:1979qx, Nester:1998mp, BeltranJimenez:2019odq}. The reason is that -- as shown -- the multiplier $\TLambda[1]{_{\mu}^{\nu\sigma}}$ changes the spectrum. The same happens in teleparallel theories for the Lagrange multiplier that enforces the vanishing of curvature (see~\cite{BeltranJimenez:2018vdo,Rigouzzo:2022yan}).

\begin{figure}[t]
  \center
  \includegraphics[width=\linewidth]{PropagandaLite.pdf}
	\caption{\label{Propaganda} 
	Multipliers are already ubiquitous in gravity. Placement of our PV theory~\cref{DC}, relative to EC~\cite{Cartan:1922,Cartan:1923,Cartan:1924,Cartan:1925,Einstein:1925, Einstein:1928,Einstein:19282} and metric GR~\cite{Einstein:1915}.
	}
\end{figure}

\section{Illustration of strong coupling alleviation} 
We shall illustrate the mechanism by which $\TLambda[1]{_{\mu}^{\nu\sigma}}$ alleviates the strong coupling problem in a simple example -- the corresponding analysis for our theory \eqref{DC}, as output of the software HiGGS \cite{Barker:2022kdk}, is presented in~\cref{FullAnalysis}. Assuming familiarity with the Dirac algorithm~\cite{Anderson:1951ta,Bergmann:1954tc,Castellani:1981us,Chen:1998ad,Yo:1999ex,Yo:2001sy} (see good pedagogical introductions~\cite{Henneaux:1992ig,Blagojevic:2002,Golovnev:2022rui}), we consider analogs of `\emph{tetrad}' $\qe{}$, `\emph{spin connection}' $\qw{}$ and `\emph{torsion}' $\dqe{}+\qw{}$ in a minimal working example (MWE) \begin{equation}\label{MWE}
	\begin{aligned}
		L_{\text{MWE}}&\equiv\qw{}\left[\qe{}\left(1+\qw{}\right)-1\right]\left(\dqw{}+1\right)\\
			&\ \ \ \ \ +\left(\qe{}-1\right)^2+\lambda\left(\dqe{}+\qw{}\right) \;,
	\end{aligned}
\end{equation}
with `\emph{Minkowski background}' $\qe{}-1\approx\qw{}\approx0$.
First try omitting the multiplier: the definitions $p_i\equiv\partial/\partial\dot{q}_i L_{\text{MWE}}\left(\lambda\mapsto 0\right)$ engender two \emph{primary} constraints $\phie{}\equiv\pie{}\approx 0$ and $\phiw{}\equiv\piw{}-\qw{}\left[\qe{}\left(1+\qw{}\right)-1\right]\approx 0$, so the Hamiltonian is $H_{\text{MWE}}\left(\lambda\mapsto 0\right)\equiv\sum_i u_i\phi_i-\qw{}\left[\qe{}\left(1+\qw{}\right)-1\right]-\left(\qe{}-1\right)^2$, where $u_i$ replace Dirac's missing/uninvertible velocities $\dot{q}_i$. The Poisson bracket $\left\{\phie{},\phiw{}\right\}\approx \qw{}\left(1+\qw{}\right)$ has the sickly property that it \emph{vanishes} in the background linearisation of~\cref{MWE}. This feature is the popular understanding of strongly coupled torsion~\cite{Cheng:1988zg,Chen:1998ad,Yo:1999ex,Yo:2001sy,Blixt:2020ekl,Barker:2021oez}. For a vanishing bracket, the consistency conditions $\dot{\phi}_i\equiv\left\{\phi_i,H_{\text{MWE}}\right\}\approx0$ engender \emph{secondary} constraints $\chie{}\equiv 2\left(\qe{}-1\right)+\qw{}\left(\qw{}+1\right)\approx 0$ and $\chiw{}\equiv\qe{}\left(1+2\qw{}\right)-1\approx 0$. Since $\{\chie{},\phie{}\}\approx 2$ and $\{\chiw{},\phiw{}\}\approx 2\qe{}$ do not vanish, so $\dot{\chi}_i\approx 0$ solve for the $u_i$ and terminate the algorithm with \emph{zero} dynamical d.o.f. In the nonlinear theory $\dot{\phi}_i\approx 0$ \emph{already} solve for the $u_i$, the $\chi_i$ are not induced, and a single \dof (two unconstrained Cauchy data) is \emph{activated}~\cite{Golovnev:2022rui}. Next, with $\lambda$ we additionally have $\phil{}\equiv\pil{}\approx 0$ and this time it is a trivial exercise to confirm that background linearisation \emph{changes nothing} (see~\cref{ConstraintAlgebra}). In this metaphor, we compare $\lambda$ in~\cref{MWE} with $\TLambda[1]{^\sigma_{\mu\nu}}$ in~\cref{DC}, the effect is to drag the nonlinear \dof down into the linear regime, and thereby solve a major problem in non-Riemannian gravity~\cite{Kopczynski:1982,Cheng:1988zg,Dimakis:1989az,Dimakis:1989ba,Lemke:1990su,Hecht:1990wn,Chen:1998ad,Yo:1999ex,Yo:2001sy,BeltranJimenez:2019hrm,Blixt:2020ekl,Barker:2021oez}.

\begin{figure}[t]
  \center
  \includegraphics[width=\linewidth]{CombinedTheories.pdf}
	\caption{\label{ConstraintAlgebra} 
	Our solution to non-Riemannian strong coupling can easily be understood through the 1D minimal working example in~\cref{MWE}. Without our method, the $\chi_e$, $\chi_\omega$ constraints convert into a nonlinear torsion d.o.f. Accordingly, in actual nonlinear gravity~\cite{Yo:2001sy,Barker:2021oez}, constraints are converted into the $\B{_{\mu\nu}}$ field in~\cref{KalbRamondEffectiveTheory}.
	}
\end{figure}

\section{Application to dark matter}


We can broadly distinguish two situations, according to whether or not the Proca field $\T[2]{_{\mu}}$ propagates at the relevant energy scales. The first case has been considered in~\cite{Bekenstein:1971hc,Ford:1989me,Armendariz-Picon:2004say,Koivisto:2007bp,Koivisto:2008xf,Hambye:2008bq,Golovnev:2008cf,Chiba:2008eh,Koivisto:2008xf,Kanno:2008gn,Hambye:2009fg,Arina:2009uq,Hisano:2010yh,Diaz-Cruz:2010czr,Lebedev:2011iq,Farzan:2012hh,Baek:2012se,Tasinato:2014eka,DeFelice:2016yws,DeFelice:2016uil,BeltranJimenez:2016afo,Belyaev:2016icc,Emami:2016ldl,Heisenberg:2017hwb,Rodriguez:2017wkg,deFelice:2017paw,Nakamura:2018oyy,Arcadi:2020jqf,Heisenberg:2020xak,Benisty:2021sul,deRham:2021efp,Barman:2021qds,Garcia-Saenz:2021uyv}. Whereas all these works used as a starting point an \emph{effective} Lagrangian with a structure similar to \cref{MasterLagrangian}, our theory \eqref{DC} can endow such models with a fundamental origin in EC gravity.


Formula \eqref{Mass2} for the vector mass hints towards the second case, in which torsion does not propagate at the relevant energy scales. Namely, we already know that $\alp{5}<0$ and $\bet{2}>-1/3$. If $\bet{2}$ is not too close to the border value (\eg $\bet{2}\gtrsim0$) and $|\alp{5}|$ is at most of order one, we observe that $\mass[2]{}^2 \gtrsim \Planck^2$. Thus, one is tempted to conclude that parameter choices for which the Proca field does not propagate below the Planck scale are more generic. This is good news since then the energy scale at which scattering amplitudes involving the massive vector field violate perturbative unitarity also lies above the Planck scale. Consequently, introducing propagating torsion with the theory \eqref{DC} generically does not lower the cutoff scale, above which the theory ceases to be predictive. 

Integrating out the field $\T[2]{_{\mu}}$ at energies below $\mass[2]{}$, we get from \eq \eqref{MasterLagrangian} (with  $\T[3]{_{\mu}}\mapsto 0$)
\begin{equation} \label{effectiveLagrangian}
	\mathcal{L}_{\text{T}} = - \frac{\Planck^2}{2}  \rR{} - \frac{1}{18 \mass[2]{}^2} \Si[2]{_{\mu}} \Si[2]{^{\mu}} + 	L_{\text{M}}\left(\rCon{}\right) \;.
\end{equation}
As a concrete example, we shall focus on fermionic fields in $L_{\text{M}}\left(\rCon{}\right)$ and first consider a generic $\Si[2]{_{\mu}} = -3 \sum_j \big(\zeta_V V^{(j)}_\mu + \zeta_A A^{(j)}_\mu\big)$, where $j$ sums species, $\zeta_V$ and $\zeta_A$ are real constants and $V_\mu = \bar{\Psi} \gamma_\mu \Psi$ and $A_\mu = \bar{\Psi} \gamma_5 \gamma_\mu \Psi$ represent the fermionic vector and axial currents, respectively. Following~\cite{Shaposhnikov:2020aen}, we shall take into account the field content of the Standard Model and add to it a singlet fermion $N$ of mass $M_N$. Such a scenario is strongly motivated since a small mixing with active neutrinos can generate neutrino masses via the seesaw mechanism (see~\cite{Boyarsky:2018tvu} for a review). For the present discussion, however, we shall not specify the precise nature of $N$. Independently of a possible connection to neutrinos, a singlet fermion is a dark matter candidate, but this option is only viable if a sufficient production of $N$ takes places in the early Universe. 

This can be achieved both by a propagating torsion field acting as mediator~\cite{Barman:2019mlj} and by the effective 4-fermion interaction arising from \eq \eqref{effectiveLagrangian} ~\cite{Shaposhnikov:2020aen}. We focus on the second case and consider the parameter choice $\zeta_V \gg 1$ and $\zeta_A \gg 1$, which can arise from a non-minimally coupled kinetic term of the fermion, $i/2 \bar{\Psi} (1+ 2 i \zeta_V + 2 i \zeta_A \gamma_5) \gamma^\mu \mathcal{D}_\mu \Psi + \text{h.c.}$ (see~\cite{Shaposhnikov:2020frq,Karananas:2021zkl,Rigouzzo:2023sbb}).  Then the produced abundance of fermions is~\cite{Shaposhnikov:2020aen}
\begin{equation}
	\frac{\Omega_N}{\Omega_{DM}}\simeq 1.7 \, C \, \frac{M_P}{\mass[2]{}} \left( \frac{M_N}{10~\text{keV}}\right) \, \left( \frac{T_\mathrm{prod}}{\mass[2]{}} \right)^3 \;.
	\label{darkMatterAbundance}
\end{equation}
Here $T_\mathrm{prod}$ is the highest temperature at which fermions are generated, \ie generically the temperature of the hot Big Bang, and $C = \left(15(\zeta_V^2 - \zeta_A^2)^2 + 7(\zeta_V + \zeta_A)^4 + 8 (\zeta_V - \zeta_A)^4\right)$ for a Dirac fermion $N$. Choosing $C$ sufficiently large, producing all of dark matter, $\Omega_N=\Omega_{DM}$, can be achieved for a wide range of masses $M_N$ down to $\text{few}~\text{keV}$, where even lighter fermions are excluded due to well-known observational bounds on warm dark matter~\cite{Garzilli:2018jqh, Garzilli:2019qki}. The resulting mass hierarchy typically is $\mass[2]{}\gg M_P \gg T_\mathrm{prod}$. Also $\mass[2]{}< M_P$ is possible, as long as $T_\mathrm{prod}<\mass[2]{}$ so that $\T[2]{_{\mu}}$ does not propagate.

\section{Conclusions} 
Several long-standing problems related to gravity have remained unsolved. Since many of the proposed solutions call for the introduction of new particles, one naturally wonders if GR can provide them and indeed it is easy to generate additional propagating degrees of freedom, \eg from higher powers of curvature terms. However, such models are generically plagued by problems due to ghost instabilities~\cite{Stelle:1977ry,Neville:1978bk,Neville:1979rb,Sezgin:1979zf,Hayashi:1979wj, Hayashi:1980qp} and strong coupling~\cite{Hecht:1996np,Yo:1999ex,Yo:2001sy}. In this paper, we have simultaneously addressed both issues by providing the first example of propagating vector torsion  arising from a curvature-squared invariant in EC gravity, for which consistency can be proven at the full nonlinear level.

Our theory \eqref{DC} endows effective Einstein-Proca models with a fundamental gravitational origin. This leads to significant constraints on the parameter space even in situations in which torsion does not propagate, and one can say that the remaining freedom is `just right'. On the one hand, it is possible to address some of the unresolved puzzles in cosmology with approaches that do not exist in the metric formulation of GR. As an example, we discussed a novel mechanism for producing fermionic dark matter~\cite{Shaposhnikov:2020aen}. On the other hand, theories such as the one that we have constructed generically feature significantly fewer free parameters as compared to effective Einstein-Proca models~\cite{Bekenstein:1971hc,Ford:1989me,Armendariz-Picon:2004say,Koivisto:2007bp,Golovnev:2008cf,Chiba:2008eh,Koivisto:2008xf,Kanno:2008gn,Koivisto:2008xf,Hambye:2008bq,Hambye:2009fg,Arina:2009uq,Hisano:2010yh,Diaz-Cruz:2010czr,Lebedev:2011iq,Farzan:2012hh,Baek:2012se,Tasinato:2014eka,DeFelice:2016yws,Emami:2016ldl,DeFelice:2016uil,BeltranJimenez:2016afo,deFelice:2017paw,Heisenberg:2017hwb,Rodriguez:2017wkg,Nakamura:2018oyy,Heisenberg:2020xak,Arcadi:2020jqf,Benisty:2021sul,deRham:2021efp,Barman:2021qds,Garcia-Saenz:2021uyv} or proposals~\cite{Freidel:2005sn,Bauer:2008zj,Poplawski:2011xf,Diakonov:2011fs,Khriplovich:2012xg,Magueijo:2012ug,Khriplovich:2013tqa,Markkanen:2017tun,Carrilho:2018ffi,Enckell:2018hmo,Rasanen:2018fom,Rubio:2019ypq,Shaposhnikov:2020geh,Karananas:2020qkp,Langvik:2020nrs,Shaposhnikov:2020gts,Mikura:2020qhc,Shaposhnikov:2020aen,Kubota:2020ehu,Enckell:2020lvn,Iosifidis:2021iuw,Racioppi:2021ynx,Cheong:2021kyc,Dioguardi:2021fmr,Piani:2022gon,Dux:2022kuk,Rasanen:2022ijc,Gialamas:2023emn,Gialamas:2023flv,Piani:2023aof,Poisson:2023tja,Rigouzzo:2023sbb} with non-propagating torsion, and hence are more predictive. 

Despite the attractive by-products in this case (\eg dark matter production), it may be argued that there is \emph{not presently a fundamental need} for a new massive vector in cosmology and fundamental physics and many of the open issues, such as inflation and dark energy, can be addressed more easily with scalar fields. It is interesting therefore that the theoretical efforts required to extract such a vector from torsion are quite strenuous. Our paper shows this very clearly (see~\cref{WhacAMole}); the torsion vectors are badly entangled with the tensor mode. So one might even conclude that the conceptual difficulties in constructing a theory with a propagating torsion vector match nicely with the absence of phenomenological indications for the existence of an additional vector field.

Several directions for future research emerge from our findings. First, our method of analysis has the potential to rule out many of the models~\cite{Sezgin:1981xs,Blagojevic:1983zz,Blagojevic:1986dm,Kuhfuss:1986rb,Yo:1999ex,Yo:2001sy,Blagojevic:2002,Puetzfeld:2004yg,Yo:2006qs,Nair:2008yh,Shie:2008ms,Chen:2009at,Nikiforova:2009qr,Ni:2009fg,Baekler:2010fr,Ho:2011qn,Ho:2011xf,Ong:2013qja,Puetzfeld:2014sja,Karananas:2014pxa,Karananas:2014pxa,Ho:2015ulu,Ni:2015poa,Karananas:2016ltn,Obukhov:2017pxa,Blagojevic:2017ssv,Blagojevic:2018dpz,Lin:2018awc,Blagojevic:2018dpz,Tseng:2018feo,Blagojevic:2018dpz,Zhang:2019xek,Zhang:2019mhd,Jimenez:2019qjc,Lin:2019ugq,BeltranJimenez:2019acz,Aoki:2019rvi,Percacci:2019hxn,BeltranJimenez:2020sqf,Barker:2020gcp,MaldonadoTorralba:2020mbh,Marzo:2021esg,delaCruzDombriz:2021nrg,Marzo:2021iok,Baldazzi:2021kaf,Barker:2021oez,Annala:2022gtl} that -- as a result of linear study only -- are seemingly consistent. Second, we have developed an approach to construct consistent theories. In particular, it remains to be determined if our theory is special or if classes of models with analogous properties exist in EC gravity and other formulations of GR. Third, our model can constrain the parameter space in inflationary models that suffer from a loss of predictivity due to numerous unknown coupling constants (see \eg \cite{Bauer:2008zj,Rodriguez:2017wkg,Rasanen:2018ihz, Raatikainen:2019qey, Langvik:2020nrs, Shaposhnikov:2020gts,Rigouzzo:2022yan}) and have further observables consequences such as in gravitational waves~\cite{Battista:2021rlh}. Finally, consistent gravitational theories with curvature-squared terms may provide new approaches for the ultimate challenge of UV-completing GR~\cite{Lin:2019ugq,Melichev:2023lwj}.

\begin{acknowledgments}
This work was performed using resources provided by the Cambridge Service for Data Driven Discovery (CSD3) operated by the University of Cambridge Research Computing Service (\href{www.csd3.cam.ac.uk}{www.csd3.cam.ac.uk}), provided by Dell EMC and Intel using Tier-2 funding from the Engineering and Physical Sciences Research Council (capital grant EP/T022159/1), and DiRAC funding from the Science and Technology Facilities Council (\href{www.dirac.ac.uk}{www.dirac.ac.uk}).

We thank Jaakko Annala, Daniel Blixt, Francisco José Maldonado Torralba, Carlo Marzo, Roberto Percacci, Syksy Räsänen, Claire Rigouzzo, Inar Timiryasov and Tom Złośnik for useful comments. Moreover we are indebted to anonymous referees for insightful feedback. W.B. is grateful for the kind hospitality of Leiden University and the Lorentz Institute, and the support of Girton College, Cambridge. S.Z.~acknowledges support of the Fonds de la Recherche Scientifique -- FNRS.
\end{acknowledgments}

\appendix

\section{Conventions}\label{Conventions}

In this appendix we quote our conventions. The metric $g_{\mu\nu}$ has signature $\left(+,-,-,-\right)$ and the covariant derivative of a covector $A_\nu$ is $\rcD{_\mu} A_\nu = \partial_\mu A_\nu - \Gamma^\alpha_{~~ \mu\nu} A_\alpha$ with an a priori independent affine connection $\Gamma^\alpha_{~~\beta \gamma}$ and torsion $\tensor{T}{^\mu_{\nu\sigma}}\equiv -2\rcCon{^\mu_{[\nu\sigma]}}\equiv\frac{4}{3}\T[1]{^\mu_{[\nu\sigma]}}+\frac{2}{3}\tensor*{\delta}{^\mu_{[\nu}}\T[2]{_{\sigma]}}+\tensor{\epsilon}{^\mu_{\nu\sigma\lambda}}\T[3]{^\lambda}$ decomposed into tensor $\T[1]{^\lambda_{\mu\nu}}\equiv\frac{1}{2}\tensor{T}{^\lambda_{\mu\nu}}+\tfrac{1}{2}\tensor{T}{_\mu^\lambda_\nu}+\tfrac{1}{6}\tensor{g}{_{\mu\nu}}\T[2]{^\lambda}-\tfrac{1}{6}\tensor*{\delta}{^\lambda_\nu}\T[2]{_\mu}+\tfrac{1}{3}\tensor*{\delta}{^\lambda_\mu}\T[2]{_\nu}$, vector $\T[2]{_\mu}\equiv\tensor{T}{^\lambda_{\lambda\mu}}$ and pseudovector $\T[3]{_\mu}\equiv\tfrac{1}{6}\tensor{\epsilon}{_{\mu\nu\sigma\lambda}}\tensor{T}{^{\nu\sigma\lambda}}$ parts, with nonmetricity $Q_{\alpha\mu\nu}\equiv\nabla_\alpha g_{\mu\nu}$, where anti-symmetrization on some tensor $T_{[\mu\nu]}$ is defined as $T_{[\mu\nu]} \equiv \frac{1}{2}(T_{\mu\nu}-T_{\nu\mu})$. The curvature has contractions $\rcR{}\equiv\rcR{^\lambda_\lambda}\equiv\rcR{^{\mu\nu}_{\mu\nu}}$, with $\rcR{_\lambda^\sigma_{\mu\nu}}\equiv 2\PD{_{[\nu}}\rcCon{^\sigma_{\mu]\lambda}}+2\rcCon{^\sigma_{[\nu|\kappa}}\rcCon{^\kappa_{|\mu]\lambda}}$, and the Riemann curvature $\rR{_\lambda^\sigma_{\mu\nu}}$ is analogously defined in terms of the metric-compatible and torsion-free Levi-Civita connection $\rCon{^\sigma_{\mu\nu}}$, with covariant derivative $\rD{_\mu}$. We also use the tetrad $\tensor{e}{_i^\mu}\tensor{e}{_j^\nu}\tensor{\eta}{^{ij}}\equiv\tensor{g}{^{\mu\nu}}$, with $e\equiv\det\tensor{e}{^i_\mu}\equiv\sqrt{-g}\equiv\sqrt{-\det\tensor{g}{_{\mu\nu}}}$, and independent spin connection $\tensor{\omega}{^{ij}_\mu}$. For our gamma matrices, we use the convention $\left\{\gamma_i, \gamma_j \right\} \equiv - 2 \eta_{ij}$, $\gamma_5 \equiv i \gamma_0 \gamma_1 \gamma_2 \gamma_3$, with Greek indices on gamma matrices implying contraction with an (inverse) tetrad field, and in flat space with Cartesian coordinates the Levi-Civita tensor has the same components as the symbol, $\epsilon_{0123}\equiv 1$.

\begin{figure}[t]
  \center
  \includegraphics[width=\linewidth]{PhaseSpace.pdf}
	\caption{\label{PhaseSpace} 
	Strong coupling splits the perturbative vacuum. Left: the linear approximation to~\cref{ToyModel} erroneously anticipates the stable vacuum. Right: in the bulk of the true nonlinear phase space this vacuum is torn apart by a high-energy separatrix.
	}
\end{figure}

\begin{figure*}[h]
  \center
	\vspace{20pt}
  \includegraphics[width=\linewidth]{SupplementalCombinedTheories.pdf}
	\caption{\label{FullAnalysisFigure} 
	On-shell constraint algebra of $\LPV{(2)\rightleftarrows(3)}$ in~\eqref{DC}, with and without our multipliers, with the logical flow of the \emph{linear} Dirac algorithm superimposed. This figure supplements~\cref{ConstraintAlgebra} for completeness. Without our method, the induced secondary constraints $\tensor*{\chi}{_e^{1^-}}$, $\tensor*{\chi}{_\omega^{0^-}}$, $\tensor*{\chi}{_\omega^{2^-}}$ of the theory~\cref{DC} would collectively lose six (canonical) \dof in nonlinear gravity~\cite{Yo:2001sy,Barker:2021oez}, activating half of the $\B{_{\mu\nu}}$ field in~\cref{KalbRamond}. If $\bet{2}\to 0$ had additionally been enforced, the whole of $\B{_{\mu\nu}}$ would have become activated. The reason we do not show the nonlinear algebra in this case is because the loss of secondaries occurs \emph{among} specific components of given spin sectors, so the convenient spin-parity decomposition of the algebra falls apart. The full nonlinear constraint chains are listed in~\cite{HamiltonianAnalysis}.
	}
\end{figure*}

\section{Toy-model of strong coupling}\label{SCIntro}

\vspace{10pt}

In this appendix, we clarify what we mean by \emph{strong coupling}, and why it is commonly diagnosed in the Hamiltonian framework. Consider the double-oscillator toy model (TM) of \dof $q_1$ and $q_2$ in one dimension
\begin{equation}\label{ToyModel}
	L_{\text{TM}}=\frac{1}{2}\dot{q}_1^2-\frac{1}{2}q_1^2+\frac{1}{2}q_1^2\dot{q}_2^2-\frac{1}{2}q_2^2 \;.
\end{equation}
The perturbative `approximation' to~\cref{ToyModel} is $L_{\text{TM}}=\dot{q}_1^2/2-q_1^2/2-q_2^2/2+\mathcal{O}\left(q^4\right)$, for which the second oscillator drops out algebraically as $q_2\approx 0$, leaving the Hamiltonian phase subspace $\left(q_1,p_1\equiv\dot{q}_1\right)$ in which the first oscillator orbits the perfectly healthy vacuum at $q_1\approx p_1\approx 0$ with Hamilton's equations $\left(\dot{q}_1\approx p_1,\dot{p}_1\approx-q_1\right)$. However, the true physics always lives in the space $\left(q_1,p_1,q_1,p_2\right)$, in which the Hamilton equation $\dot{p}_1\approx -q_1+p_2^2/q_1^2$ is \emph{inherently non-perturbative}. As illustrated in~\cref{PhaseSpace}, the effect is absolutely fatal to the perturbative vacuum. Generalising to field theory, strongly-coupled \dof, analogous to $q_2$, reveal themselves via structural changes to the Hamiltonian constraint structure when passing from linearised to nonlinear gravity theories. In such cases it is not clear how the background of the linearisation can then be a viable spacetime, even though it may perfectly well be an exact solution to the nonlinear field equations. The sick kinetic structure of~\eqref{ToyModel} may be compared to that in~\eqref{KalbRamondEffectiveTheory}.

It is important to point out that the splitting of the vacuum in~\cref{PhaseSpace} is an effect which \emph{need not} be associated with ghost modes. Indeed, in~\cref{ToyModel} we are careful to preserve the sign of the kinetic term. In our practical, field theoretic example in~\cref{KalbRamondEffectiveTheory} however, it may turn out to be harder to bound the Hamiltonian due to the multiple components of $\B{_{\mu\nu}}$. 

\section{Full nonlinear Hamiltonian analysis}\label{FullAnalysis}

In this appendix, we quote the full, nonlinear Hamiltonian analysis of $\LPV{(2)\rightleftarrows(3)}$ corresponding to $\alp{5}>0$ w.l.o.g. The full Dirac algorithm is shown schematically in~\cref{FullAnalysisFigure}, using precisely the same notation as in~\cref{ConstraintAlgebra}. Note the appearance of tertiary $\zeta$ and quaternary $\theta$ constraints. The full nonlinear constraint chains are found in~\cite{HamiltonianAnalysis}, these are fully computed by the HiGGS software based only on the definition of the Lagrangian~\cite{Barker:2022kdk}. The key observation is that the size of the algebra (i.e. the number of induced non-primary constraints) does not change in passing from the linear to the nonlinear theories (though we do not exclude this effect at other points which punctuate the bulk of the phase space). This is consistent with the \emph{manifestly} healthy and dynamically equivalent Einstein-Proca model.

\bibliography{Manuscript}

\end{document}
